\documentclass[11pt]{article}

\usepackage{amsmath}
\usepackage{amssymb}
\usepackage{amsthm}
\usepackage{mathtools}
\usepackage{stmaryrd}
\usepackage[margin=1in]{geometry}

\newtheorem{theorem}{Theorem}
\newtheorem{lemma}[theorem]{Lemma}
\newtheorem{definition}[theorem]{Definition}
\newtheorem{corollary}[theorem]{Corollary}

\newcommand{\sem}[1]{\llbracket #1 \rrbracket}
\newcommand{\diff}{\mathsf{diff}}
\newcommand{\SType}{\mathsf{SType}}
\newcommand{\Tags}{\mathsf{Tags}}
\newcommand{\TType}{\mathsf{TType}}
\newcommand{\Clauses}{\mathsf{Clauses}}
\newcommand{\Clause}{\mathsf{Clause}}
\newcommand{\STypes}{\mathsf{STypes}}
\newcommand{\TInst}{\mathsf{TInst}}
\newcommand{\TInsts}{\mathsf{TInsts}}
\newcommand{\TParam}{\mathsf{TParam}}
\newcommand{\vars}{\mathsf{vars}}
\newcommand{\stypes}{\mathsf{stypes}}
\newcommand{\insts}{\mathsf{insts}}
\newcommand{\clauses}{\mathsf{clauses}}
\newcommand{\tagfn}{\mathsf{tag}}
\newcommand{\tags}{\mathsf{tags}}
\newcommand{\rtag}{\mathsf{rtag}}
\newcommand{\Val}{\mathsf{Val}}
\newcommand{\Payload}{\mathsf{Payload}}

\title{Soundness of Type Difference in SKJS}
\author{}
\date{}

\begin{document}

\maketitle

\begin{abstract}
We prove that the type difference ($\diff$) operation in SKJS satisfies 
properties ensuring sound type inference.
\end{abstract}

\section{Definitions Index}

\renewcommand{\arraystretch}{1.15}
\begin{center}
\begin{tabular}{ll}
\textbf{Symbol} & \textbf{Defined In} \\
\hline
$\Tags$ & Definition~\ref{def:runtime-tags} \\
$\Payload, \Val, \rtag$ & Definition~\ref{def:runtime-values-tags} \\
$\SType$ & Definition~\ref{def:simple-types} \\
$\tagfn$ & Definition~\ref{def:simple-types-and-tags} \\
\textbf{A0} & Definition~\ref{ax:A0} \\
\textbf{A1} & Definition~\ref{ax:A1} \\
$\STypes, \tags(\cdot)$ & Definition~\ref{def:stypes} \\
$\mathsf{STC}(\cdot,\cdot)$ & Definition~\ref{def:stc} \\
$\TInst$ & Definition~\ref{def:type-instances} \\
$\TInsts$ & Definition~\ref{def:type-instance-sets} \\
$\TParam$ & Definition~\ref{def:type-parameters} \\
$\Clause$ & Definition~\ref{def:clause} \\
$\Clauses$ & Definition~\ref{def:clauses} \\
$\TType$ & Definition~\ref{def:ttype} \\
$\diff$ (soundness P1/P2) & Definition~\ref{def:diff-soundness} \\
\end{tabular}
\end{center}

\section{Type Representation}

\begin{definition}[Runtime Tags]
\label{def:runtime-tags}
Let $\Tags$ be the set of runtime tags (the coarse runtime kind of a value). In the current implementation, tags are represented as integers, so one may take $\Tags = \{0,1,\ldots,11\}$.
\end{definition}

\begin{definition}[Runtime Values and Tags]
\label{def:runtime-values-tags}
Assume runtime values carry their tag as part of the representation: let $\Payload$ be an abstract set of payloads and define
\[
\Val = \Tags \times \Payload.
\]
Define the runtime tag function $\rtag : \Val \to \Tags$ as the first projection:
\[
\rtag((k,p)) = k.
\]
\end{definition}

\begin{definition}[Simple Types]
\label{def:simple-types}
Let $\SType$ be the set of \emph{simple types}, i.e. the atomic type forms that have a single runtime tag.
Formally, $\SType$ is the smallest set generated by the following constructors:
\begin{itemize}
\item primitive atoms: \texttt{undefined}, \texttt{void}, \texttt{null}, \texttt{boolean}, \texttt{number}, \texttt{string}, \texttt{bigint}
\item boolean literals: \texttt{true}, \texttt{false}
\item numeric literals: $n$ for $n \in \mathbb{R}$
\item string literals: \text{``}x\text{''} for strings $x$
\item structured atoms: \texttt{tuple}$(\ldots)$, \texttt{Promise}$(\ldots)$, \texttt{object}$(\ldots)$, \texttt{algebraic}$(\ldots)$ (tagged unions/ADTs with a runtime discriminant), \texttt{function}$(\ldots)$
\end{itemize}
We write these constructors suggestively; their internal parameters are left abstract, as they are not needed for the tag-level reasoning in this note.
\end{definition}

\begin{definition}[Simple Types and Tags]
\label{def:simple-types-and-tags}
The function $\tagfn : \SType \to \Tags$ assigns tags:
\begin{align*}
\tagfn(\texttt{undefined}) &= 0 & \tagfn(\texttt{void}) &= 1 \\
\tagfn(\texttt{null}) &= 2 & \tagfn(\texttt{boolean}), \tagfn(\texttt{true}), \tagfn(\texttt{false}) &= 3 \\
\tagfn(\texttt{number}), \tagfn(n) &= 4 \text{ for } n \in \mathbb{R} & 
\tagfn(\texttt{string}), \tagfn(\text{``}x\text{''}) &= 5 \\
\tagfn(\texttt{bigint}) &= 6 & \tagfn(\texttt{tuple}(\ldots)) &= 7 \\
\tagfn(\texttt{Promise}(\ldots)) &= 8 & \tagfn(\texttt{object}(\ldots)) &= 9 \\
\tagfn(\texttt{algebraic}(\ldots)) &= 10 & \tagfn(\texttt{function}(\ldots)) &= 11
\end{align*}
The semantic interpretation $\sem{\cdot} : \SType \to \mathcal{P}(\Val)$ 
maps each simple type to its set of inhabitants.
\paragraph{Note.}
We do \emph{not} define $\sem{\cdot}$ in this document; it is treated as an abstract interpretation function.
Instead, we state explicit assumptions about its behavior as axioms (A0--A1).
\end{definition}

\begin{definition}[\textbf{A0} (Runtime Tag Consistency)]
\label{ax:A0}
The semantic interpretation of simple types agrees with runtime tags:
\[
\forall s \in \SType.\; \forall v \in \Val.\; v \in \sem{s} \implies \rtag(v) = \tagfn(s).
\]
\end{definition}

\begin{definition}[\textbf{A1} (Tag Disjointness)]
\label{ax:A1}
\[
\forall s_1, s_2 \in \SType.\; \tagfn(s_1) \neq \tagfn(s_2) \implies \sem{s_1} \cap \sem{s_2} = \emptyset
\]
\end{definition}

\begin{definition}[STypes]
\label{def:stypes}
$\STypes$ is a finite partial map $\Tags \rightharpoonup \SType$ 
with at most one entry per tag: if $S(k) = s$ then $\tagfn(s) = k$.
\[
\sem{S} = \bigcup_{k \in \mathrm{dom}(S)} \sem{S(k)}
\]
We write $\tags(S) = \mathrm{dom}(S)$.
We write $\{k \mapsto s\}$ for the singleton partial map with domain $\{k\}$.
\paragraph{Intuition.}
An $\STypes$ value represents the \emph{tag-level} part of a union type: it records, for each runtime tag $k$ that may occur, a single representative simple type $s$ of that tag.
The map shape enforces a normalization invariant (``at most one entry per tag''), which matches the implementation: many operations reason by tags and intentionally do not keep multiple alternatives of the same tag.
\paragraph{Example.}
The union \texttt{number | string} is represented as:
\[
S = \{4 \mapsto \texttt{number},\; 5 \mapsto \texttt{string}\}
\]
and has semantics $\sem{S} = \sem{\texttt{number}} \cup \sem{\texttt{string}}$.
In contrast, \texttt{number} and the numeric literal \texttt{1} share the same tag $4$, so an $\STypes$ value cannot simultaneously store both; this is why, for example, a tag-based ``diff'' cannot precisely express \texttt{number \textbackslash{} 1}.
\end{definition}

\begin{definition}[Type Instances]
\label{def:type-instances}
Let $\TInst$ be the set of type instances (type variables or generic instantiations).
\paragraph{Examples.}
A type variable instance is something like $T$ (a unification variable introduced during type checking; this is \emph{not} a source-level type parameter).
A generic instantiation instance is a named type applied to type arguments, e.g. $\texttt{Promise}\langle \texttt{number} \rangle$ or $\texttt{Array}\langle \texttt{string} \rangle$.

Each instance $i \in \TInst$ has a semantic interpretation $\sem{i} \subseteq \Val$.
\end{definition}

\begin{definition}[Type Instance Sets]
\label{def:type-instance-sets}
Let $\TInsts = \mathcal{P}_{\mathrm{fin}}(\TInst)$ be the set of \emph{finite} sets of instances.
For $I \in \TInsts$:
\[
\sem{I} = \bigcup_{i \in I} \sem{i}.
\]
\end{definition}

\begin{definition}[Type Parameters]
\label{def:type-parameters}
Let $\TParam$ be the set of \emph{type parameters} (generic binders), i.e. names introduced by source-level generic declarations such as $\texttt{function f}\langle T\rangle(\ldots)$ or $\texttt{type Box}\langle T\rangle = \ldots$.
They differ from type-variable instances: type parameters are binders whose meaning is fixed by the caller/instantiation, whereas type-variable instances are inference unknowns whose meaning is solved by the type checker.
\end{definition}

\begin{definition}[Clause]
\label{def:clause}
A clause $c = (\vars, \stypes)$ consists of:
\begin{itemize}
\item $\vars \subseteq \TParam$: a finite set of type parameters
\item $\stypes \in \{\mathsf{None}\} \cup \{\mathsf{Some}(S) \mid S \in \STypes\}$
\end{itemize}
Each type parameter $v \in \TParam$ has an interpretation $\sem{v} \subseteq \Val$.
Semantics:
\[
\sem{(\vars, \stypes)} = \left(\bigcap_{v \in \vars} \sem{v}\right) \cap \sem{\stypes}
\]
where $\sem{\mathsf{None}} = \Val$ (top) and $\sem{\mathsf{Some}(S)} = \sem{S}$.
For $\vars = \emptyset$, $\bigcap_{v \in \emptyset} \sem{v} = \Val$.
\[
\text{Well-formedness: } \stypes = \mathsf{None} \implies \vars = \emptyset \quad (\text{i.e., } (\emptyset,\mathsf{None}) \text{ is the unique top clause}).
\]
\end{definition}

\begin{definition}[Clauses]
\label{def:clauses}
$\Clauses$ is a finite map from $\mathcal{P}_{\mathrm{fin}}(\TParam)$ to $\Clause$, 
with at most one clause per key, and such that keys agree with clause variables: if $C(k) = (\vars,\stypes)$ then $\vars = k$.
\[
\sem{C} = \bigcup_{k \in \mathrm{dom}(C)} \sem{C(k)}
\]
\paragraph{Example.}
Let $S_{\mathsf{num}} = \{4 \mapsto \texttt{number}\}$ and $S_{\mathsf{str}} = \{5 \mapsto \texttt{string}\}$.
Let $A,B \in \TParam$ be two type parameters.
Consider the clauses map $C$ with two entries:
\[
C(\{A\}) = (\{A\}, \mathsf{Some}(S_{\mathsf{num}}))
\qquad
C(\{B\}) = (\{B\}, \mathsf{Some}(S_{\mathsf{str}})).
\]
Then
\[
\sem{C} =
\bigl(\sem{A} \cap \sem{\texttt{number}}\bigr)
\cup
\bigl(\sem{B} \cap \sem{\texttt{string}}\bigr).
\]
To illustrate a non-singleton key, consider instead the single-entry map $C'$ defined by:
\[
C'(\{A,B\}) = (\{A,B\}, \mathsf{Some}(S_{\mathsf{num}})).
\]
Then $\sem{C'} = \sem{A} \cap \sem{B} \cap \sem{\texttt{number}}$, i.e. the clause requires \emph{both} parameters.
\end{definition}

\begin{definition}[TType]
\label{def:ttype}
$\TType = (\insts, \clauses)$ where $\insts \in \TInsts$ and $\clauses \in \Clauses$.
\[
\sem{(\insts, \clauses)} = \sem{\insts} \cup \sem{\clauses}
\]
\paragraph{Example.}
Let $I = \{\texttt{Promise}\langle \texttt{number} \rangle\} \in \TInsts$ and let $C$ be the clauses map from the previous example.
Then $t = (I,C)$ denotes the type whose semantics is:
\[
\sem{t} = \sem{\texttt{Promise}\langle \texttt{number} \rangle} \cup \sem{C}.
\]
\end{definition}

\section{The Diff Operation}

\begin{definition}[Diff Soundness Properties]
\label{def:diff-soundness}
Let $X$ be a set equipped with a semantic interpretation $\sem{\cdot} : X \to \mathcal{P}(\Val)$.
An operation $\diff : X \times X \to X$ is \emph{sound} (with respect to $\sem{\cdot}$) if for all $x_1,x_2 \in X$, letting $x_3 = \diff(x_1,x_2)$, the following properties hold:
\begin{align}
\text{(P1)} \quad & \sem{x_3} \subseteq \sem{x_1} \tag{under-approximation} \\
\text{(P2)} \quad & \sem{x_3} \cup \sem{x_2} \supseteq \sem{x_1} \tag{coverage}
\end{align}
In this document we instantiate $X$ with $\STypes$, $\Clause$, $\Clauses$, and $\TType$.
\end{definition}

\begin{lemma}[Equivalent Characterization]
(P1) $\land$ (P2) $\iff$ $\sem{t_1} \setminus \sem{t_2} \subseteq \sem{t_3} \subseteq \sem{t_1}$.
\end{lemma}

\begin{proof}
($\Rightarrow$) Let $x \in \sem{t_1} \setminus \sem{t_2}$. By (P2), 
$x \in \sem{t_3} \cup \sem{t_2}$. Since $x \notin \sem{t_2}$, $x \in \sem{t_3}$.
Thus $\sem{t_1} \setminus \sem{t_2} \subseteq \sem{t_3}$. Combined with (P1): 
$\sem{t_1} \setminus \sem{t_2} \subseteq \sem{t_3} \subseteq \sem{t_1}$.

($\Leftarrow$) (P1) is immediate. For (P2): 
$\sem{t_3} \cup \sem{t_2} \supseteq (\sem{t_1} \setminus \sem{t_2}) \cup \sem{t_2} = \sem{t_1} \cup \sem{t_2} \supseteq \sem{t_1}$.
\end{proof}

\begin{definition}[Completeness]
A diff operation is \emph{complete} if:
\[
\sem{\diff(t_1, t_2)} = \sem{t_1} \setminus \sem{t_2}
\]
By the equivalent characterization, (P1) $\land$ (P2) implies 
$\sem{t_1} \setminus \sem{t_2} \subseteq \sem{\diff(t_1, t_2)}$.
Completeness additionally requires the reverse: $\sem{\diff(t_1, t_2)} \subseteq \sem{t_1} \setminus \sem{t_2}$.
\end{definition}

\subsection{STypes Diff}

\begin{definition}[STypes Diff]
\[
\diff_{\STypes}(L, R)(k) =
\begin{cases}
L(k) & \text{if } k \in \tags(L) \setminus \tags(R) \\
\text{undefined} & \text{otherwise}
\end{cases}
\]
\end{definition}

\begin{lemma}[STypes Diff Satisfies (P1)]
$\sem{\diff_{\STypes}(L, R)} \subseteq \sem{L}$.
\end{lemma}

\begin{proof}
By definition, $\tags(\diff_{\STypes}(L,R)) = \tags(L) \setminus \tags(R) \subseteq \tags(L)$, and for all $k \in \tags(\diff_{\STypes}(L,R))$ we have $\diff_{\STypes}(L,R)(k)=L(k)$.
Therefore:
\[
\sem{\diff_{\STypes}(L,R)}
= \bigcup_{k \in \tags(\diff_{\STypes}(L,R))} \sem{L(k)}
\subseteq \bigcup_{k \in \tags(L)} \sem{L(k)}
= \sem{L}.
\]
\end{proof}

\begin{definition}[Same-Tag Coverage Condition]
\label{def:stc}
For $L, R \in \STypes$, define:
\[
\mathsf{STC}(L, R) \iff \forall k \in \tags(L) \cap \tags(R).\; \sem{L(k)} \subseteq \sem{R(k)}
\]
\end{definition}

\begin{lemma}[STypes Diff Satisfies (P2) under STC]
If $\mathsf{STC}(L, R)$ holds, then $\sem{\diff_{\STypes}(L, R)} \cup \sem{R} \supseteq \sem{L}$.
\end{lemma}

\begin{proof}
Let $x \in \sem{L}$. Then $x \in \sem{L(k)}$ for some $k \in \tags(L)$.
\begin{itemize}
\item Case $k \notin \tags(R)$: Then $k \in \tags(\diff_{\STypes}(L,R))$ and $\diff_{\STypes}(L,R)(k)=L(k)$, so $x \in \sem{\diff_{\STypes}(L,R)}$.
\item Case $k \in \tags(R)$: By $\mathsf{STC}$, $\sem{L(k)} \subseteq \sem{R(k)} \subseteq \sem{R}$, so $x \in \sem{R}$.
\end{itemize}
\end{proof}

\begin{corollary}[STypes Diff with Tag Disjointness]
If $\mathsf{STC}(L, R)$ holds, then (P2) holds.
In particular, if $L(k) = R(k)$ for all $k \in \tags(L) \cap \tags(R)$, then (P2) holds.
\end{corollary}

\begin{lemma}[STypes Diff Completeness]
$\diff_{\STypes}$ is complete (i.e., $\sem{\diff_{\STypes}(L, R)} = \sem{L} \setminus \sem{R}$) 
if and only if $\mathsf{STC}(L, R)$.
\end{lemma}

\begin{proof}
($\Rightarrow$) Suppose completeness holds and let $k \in \tags(L) \cap \tags(R)$.
Since $k \notin \tags(\diff_{\STypes}(L,R))$, the component $L(k)$ contributes no value to $\sem{L} \setminus \sem{R}$ under completeness, hence $\sem{L(k)} \subseteq \sem{R}$.
By A1 and the ``at most one entry per tag'' invariant for $\STypes$, $\sem{L(k)}$ can only overlap with $\sem{R(k)}$, so $\sem{L(k)} \subseteq \sem{R(k)}$.

($\Leftarrow$) We show $\sem{\diff_{\STypes}(L, R)} = \sem{L} \setminus \sem{R}$.

($\subseteq$) Let $x \in \sem{\diff_{\STypes}(L, R)}$. Then $x \in \sem{L(k)}$ for some $k \in \tags(L) \setminus \tags(R)$.
By A1, $x \notin \sem{R}$ (no type in $R$ has tag $k$), so $x \in \sem{L} \setminus \sem{R}$.

($\supseteq$) Let $x \in \sem{L} \setminus \sem{R}$. Then $x \in \sem{L(k)}$ for some $k \in \tags(L)$ and $x \notin \sem{R}$.
If $k \in \tags(R)$, then by hypothesis $\sem{L(k)} \subseteq \sem{R(k)} \subseteq \sem{R}$, contradicting $x \notin \sem{R}$.
So $k \notin \tags(R)$, meaning $k \in \tags(\diff_{\STypes}(L,R))$ and $x \in \sem{\diff_{\STypes}(L, R)}$.
\end{proof}

\begin{corollary}[When STypes Diff is Incomplete]
$\diff_{\STypes}(L, R)$ is incomplete when there exists $k \in \tags(L) \cap \tags(R)$ with $\sem{L(k)} \not\subseteq \sem{R(k)}$.
Example: $\diff_{\STypes}(\{4 \mapsto \texttt{number}\}, \{4 \mapsto \texttt{1}\}) = \emptyset$, 
but $\sem{\texttt{number}} \setminus \sem{\texttt{1}} \neq \emptyset$.
\end{corollary}

\subsection{TInsts Diff}

\begin{definition}[TInsts Diff]
\[
\diff_{\TInsts}(L, R) = L \setminus R = \{ i \in L \mid i \notin R \}
\]
\paragraph{Note.}
In the implementation, $\TInsts$ is represented as a collection of instance AST nodes and $\diff_{\TInsts}$ removes elements based on \emph{structural equality} of instances.
In this paper, that corresponds exactly to the usual mathematical set difference on the set $\TInst$ (i.e. membership/equality is the syntactic identity of instances, not semantic overlap).
\end{definition}

\begin{lemma}[TInsts Diff Satisfies (P1) and (P2)]
$\sem{\diff_{\TInsts}(L, R)} \subseteq \sem{L}$ and 
$\sem{\diff_{\TInsts}(L, R)} \cup \sem{R} \supseteq \sem{L}$.
\end{lemma}

\begin{proof}
(P1): $\diff_{\TInsts}(L, R) \subseteq L$, so $\sem{\diff_{\TInsts}(L, R)} \subseteq \sem{L}$.

(P2): Let $x \in \sem{L}$. Then $x \in \sem{i}$ for some $i \in L$.
\begin{itemize}
\item If $i \in R$: Then $x \in \sem{i} \subseteq \sem{R}$.
\item If $i \notin R$: Then $i \in \diff_{\TInsts}(L, R)$, so $x \in \sem{\diff_{\TInsts}(L, R)}$.
\end{itemize}
\end{proof}

\begin{lemma}[TInsts Diff Completeness]
$\diff_{\TInsts}$ is complete if and only if for all $i \in L \setminus R$ and $j \in R$:
$\sem{i} \cap \sem{j} = \emptyset$.
\end{lemma}

\begin{proof}
($\Rightarrow$) Suppose completeness and let $i \in L \setminus R$ and $j \in R$.
If $x \in \sem{i} \cap \sem{j}$, then $i \in \diff_{\TInsts}(L, R)$, so $x \in \sem{\diff_{\TInsts}(L, R)} = \sem{L} \setminus \sem{R}$.
But $x \in \sem{j} \subseteq \sem{R}$, contradiction.

($\Leftarrow$) We show $\sem{\diff_{\TInsts}(L, R)} = \sem{L} \setminus \sem{R}$.

($\subseteq$) Let $x \in \sem{\diff_{\TInsts}(L, R)}$. Then $x \in \sem{i}$ for some $i \in L \setminus R$.
For any $j \in R$, $\sem{i} \cap \sem{j} = \emptyset$ by hypothesis.
Thus $x \notin \sem{R}$, so $x \in \sem{L} \setminus \sem{R}$.

($\supseteq$) Let $x \in \sem{L} \setminus \sem{R}$. Then $x \in \sem{i}$ for some $i \in L$.
If $i \in R$, then $x \in \sem{R}$, contradiction. So $i \in L \setminus R = \diff_{\TInsts}(L, R)$.
\end{proof}

\subsection{Clause Diff}

\begin{definition}[Clause Diff]
For clauses $c_1 = (\vars_1, \stypes_1)$ and $c_2 = (\vars_2, \stypes_2)$:
\[
\diff_{\Clause}(c_1, c_2) = 
\begin{cases}
\bot & \text{if } \stypes_2 = \mathsf{None} \\
c_1 & \text{if } \stypes_1 = \mathsf{None} \\
(\vars_1, \mathsf{Some}(\diff_{\STypes}(S_1, S_2))) & \text{if } \stypes_1 = \mathsf{Some}(S_1), \stypes_2 = \mathsf{Some}(S_2)
\end{cases}
\]
where $\bot$ represents the empty type (dropped from results) with $\sem{\bot} = \emptyset$.
\end{definition}

\begin{lemma}[Clause Diff Satisfies (P1)]
$\sem{\diff_{\Clause}(c_1, c_2)} \subseteq \sem{c_1}$.
\end{lemma}

\begin{proof}
Case analysis:
\begin{itemize}
\item $\stypes_2 = \mathsf{None}$: Result is $\bot$. $\sem{\bot} = \emptyset \subseteq \sem{c_1}$.
\item $\stypes_1 = \mathsf{None}$: Result is $c_1$. $\sem{c_1} \subseteq \sem{c_1}$.
\item Both $\mathsf{Some}$: 
$\sem{(\vars_1, \mathsf{Some}(\diff_{\STypes}(S_1, S_2)))} = 
\left(\bigcap_{v \in \vars_1} \sem{v}\right) \cap \sem{\diff_{\STypes}(S_1, S_2)}
\subseteq \left(\bigcap_{v \in \vars_1} \sem{v}\right) \cap \sem{S_1} = \sem{c_1}$.
\end{itemize}
\end{proof}

\begin{lemma}[Clause Diff Satisfies (P2) Conditionally]
$\diff_{\Clause}$ satisfies (P2) when:
\begin{itemize}
\item Case $\stypes_2 = \mathsf{None}$: Always for well-formed clauses (since then $c_2$ is the unique top clause and $\sem{c_2} = \Val \supseteq \sem{c_1}$).
\item Case $\stypes_1 = \mathsf{None}$: Always (since result is $c_1$).
\item Case both $\mathsf{Some}$: When $\diff_{\STypes}(S_1, S_2)$ satisfies (P2).
\end{itemize}
\end{lemma}

\subsection{Clauses Diff}

\begin{definition}[Clauses Diff]
For $C_1, C_2 \in \Clauses$, let $c_{\emptyset} = C_2(\emptyset)$ if $\emptyset \in \mathrm{dom}(C_2)$, else undefined.
\[
\diff_{\Clauses}(C_1, C_2)(k) = 
\begin{cases}
\diff_{\Clause}(\diff_{\Clause}(C_1(k), C_2(k)), c_{\emptyset}) & \text{if } k \in \mathrm{dom}(C_2) \land c_{\emptyset} \text{ defined} \\
\diff_{\Clause}(C_1(k), C_2(k)) & \text{if } k \in \mathrm{dom}(C_2) \land c_{\emptyset} \text{ undefined} \\
\diff_{\Clause}(C_1(k), c_{\emptyset}) & \text{if } k \notin \mathrm{dom}(C_2) \land c_{\emptyset} \text{ defined} \\
C_1(k) & \text{if } k \notin \mathrm{dom}(C_2) \land c_{\emptyset} \text{ undefined}
\end{cases}
\]
for $k \in \mathrm{dom}(C_1)$. Clauses mapped to $\bot$ are dropped.
\paragraph{Note.}
Whenever $k \in \mathrm{dom}(C_2)$, the call $\diff_{\Clause}(C_1(k), C_2(k))$ compares clauses stored under the \emph{same key} $k$, so their variable-sets match by the $\Clauses$ well-formedness condition (keys agree with clause variables).
The only subtraction that may use a different key is via $c_{\emptyset}=C_2(\emptyset)$, which is deliberately applied to every $C_1(k)$; this corresponds to the implementation behavior and explains why the right-hand clause's variables are not used by $\diff_{\Clause}$.
\end{definition}

\begin{lemma}[Clauses Diff Satisfies (P1)]
$\sem{\diff_{\Clauses}(C_1, C_2)} \subseteq \sem{C_1}$.
\end{lemma}

\begin{proof}
Each output clause is derived from $C_1(k)$ via $\diff_{\Clause}$, which satisfies (P1).
\end{proof}

\subsection{TType Diff}

\begin{definition}[TType Diff]
For $t_1 = (\insts_1, \clauses_1)$ and $t_2 = (\insts_2, \clauses_2)$:
\[
\diff_{\TType}(t_1, t_2) = (\diff_{\TInsts}(\insts_1, \insts_2), \diff_{\Clauses}(\clauses_1, \clauses_2))
\]
\end{definition}

\begin{theorem}[TType Diff Satisfies (P1)]
$\sem{\diff_{\TType}(t_1, t_2)} \subseteq \sem{t_1}$.
\end{theorem}

\begin{proof}
\begin{align*}
\sem{\diff_{\TType}(t_1, t_2)} &= \sem{\diff_{\TInsts}(\insts_1, \insts_2)} \cup \sem{\diff_{\Clauses}(\clauses_1, \clauses_2)} \\
&\subseteq \sem{\insts_1} \cup \sem{\clauses_1} = \sem{t_1}
\end{align*}
by Lemmas 4 and 8.
\end{proof}

\begin{theorem}[TType Diff Satisfies (P2) Conditionally]
$\diff_{\TType}$ satisfies (P2) when all underlying $\diff_{\STypes}$ calls satisfy (P2).
In particular, (P2) holds when $\mathsf{STC}$ holds for all same-tag eliminations.
\end{theorem}

\begin{theorem}[TType Diff Completeness]
$\diff_{\TType}$ is complete when:
\begin{enumerate}
\item $\diff_{\TInsts}$ is complete (structurally distinct instances are semantically disjoint), and
\item $\diff_{\STypes}$ is complete for all clause pairs (same-tag implies $\sem{s} \subseteq \sem{s'}$).
\end{enumerate}
\end{theorem}

\begin{corollary}[Explicit Completeness Conditions for $\diff_{\TType}$]
Let $t_1=(\insts_1,\clauses_1)$ and $t_2=(\insts_2,\clauses_2)$.
Then $\diff_{\TType}(t_1,t_2)$ is complete (i.e. $\sem{\diff_{\TType}(t_1,t_2)}=\sem{t_1}\setminus\sem{t_2}$) under the following explicit hypotheses:
\begin{enumerate}
\item \textbf{Instance disjointness across sides:} for all $i\in\insts_1$ and $j\in\insts_2$ with $i\neq j$, we have $\sem{i}\cap\sem{j}=\emptyset$.
\item \textbf{Top-clause alignment:} if $\clauses_1$ is top (i.e. it contains the unique top clause $(\emptyset,\mathsf{None})$), then $\clauses_2$ is also top.
\item \textbf{Same-tag coverage for every clause subtraction:}
  \begin{enumerate}
  \item For every $k \in \mathrm{dom}(\clauses_1)\cap \mathrm{dom}(\clauses_2)$, if $\clauses_1(k)=(k,\mathsf{Some}(S_1))$ and $\clauses_2(k)=(k,\mathsf{Some}(S_2))$, then $\mathsf{STC}(S_1,S_2)$.
  \item If $\emptyset\in \mathrm{dom}(\clauses_2)$ and $\clauses_2(\emptyset)=(\emptyset,\mathsf{Some}(S_0))$, then for every $k\in \mathrm{dom}(\clauses_1)$ with $\clauses_1(k)=(k,\mathsf{Some}(S_1))$, we have $\mathsf{STC}(S_1,S_0)$.
  \end{enumerate}
\end{enumerate}
\end{corollary}

\section{Application to Type Inference}

\begin{definition}[Subtyping Check with Validation]
When checking $u_1 \leq u_2$ where $u_2$ contains type variable $T$:
\begin{enumerate}
\item Compute $\diff(u_1, u_2)$ and grow $T$'s lower bound.
\item Expand $u_1, u_2$ using the substitution.
\item Validate via $\mathsf{isSub}(\mathit{expand}(u_1), \mathit{expand}(u_2))$.
\end{enumerate}
\end{definition}

\begin{theorem}[Inference Soundness]
If $\mathsf{isSub}$ correctly validates subtyping on expanded types, then 
type inference is sound regardless of whether $\diff$ satisfies (P2).
\end{theorem}

\begin{proof}
The $\diff$ result only affects inference hints (bound accumulation). 
The final subtyping decision is made by $\mathsf{isSub}$ on expanded types, 
which does not depend on $\diff$ correctness.
\end{proof}

\section{Examples of Incompleteness}

We present two concrete examples where $\diff$ is incomplete, i.e., 
$\sem{\diff(t_1, t_2)} \neq \sem{t_1} \setminus \sem{t_2}$ (it may be a strict subset or strict superset).

\subsection{Example 1: Same-Tag Literal Types}

\textbf{TypeScript code:}
\begin{verbatim}
function f(x: 1 | 2): 2 {
  if (x === 1) { return 2; }
  return x;  // x should be narrowed to 2
}
\end{verbatim}

\textbf{Mathematical setup:}
Both \texttt{1} and \texttt{2} have the same runtime tag $4$.
Because $\STypes$ stores at most one $\SType$ per tag, it cannot precisely represent unions of same-tag literals such as \texttt{1 | 2}.
In practice, the implementation merges same-tag literals into a single representative $\SType$ for the tag (often a wider type such as \texttt{number}), which loses the ability to subtract a single literal while keeping another.

\textbf{Incompleteness:}
Since $\tagfn(\texttt{1}) = \tagfn(\texttt{2}) = 4$ and $\STypes$ stores at most one entry per tag,
$\diff_{\STypes}$ cannot distinguish between literals of the same tag.
For a simpler illustration:
\[
\diff_{\STypes}(\{4 \mapsto \texttt{number}\}, \{4 \mapsto \texttt{1}\}) = \emptyset
\]
But $\sem{\texttt{number}} \setminus \sem{\texttt{1}} = \text{all numbers except 1} \neq \emptyset$.

The condition $\mathsf{STC}(L, R)$ fails: $\sem{\texttt{number}} \not\subseteq \sem{\texttt{1}}$.

\subsection{Example 2: Top Type (unknown)}

\textbf{TypeScript code (requires \texttt{--strictNullChecks}):}
\begin{verbatim}
function f(x: unknown): {} {
  if (x !== null && x !== undefined) {
    return x;  // x should be narrowed to {}
  }
  return {};
}
\end{verbatim}

\textbf{Mathematical setup:}
The type \texttt{unknown} is represented as a clause with $\stypes = \mathsf{None}$ (top).
Let $c_1 = (\emptyset, \mathsf{None})$ represent \texttt{unknown}.
Let $c_2$ represent \texttt{null | undefined}: $c_2 = (\emptyset, \mathsf{Some}(\{0 \mapsto \texttt{undefined}, 2 \mapsto \texttt{null}\}))$.

\textbf{Incompleteness:}
By definition of $\diff_{\Clause}$:
\[
\diff_{\Clause}(c_1, c_2) = c_1 \quad \text{(since } \stypes_1 = \mathsf{None}\text{)}
\]
So $\sem{\diff_{\Clause}(c_1, c_2)} = \sem{c_1} = \Val$ (all values).

But the true set difference is:
\[
\sem{c_1} \setminus \sem{c_2} = \Val \setminus (\sem{\texttt{null}} \cup \sem{\texttt{undefined}}) \neq \Val
\]

The result is incomplete: $\diff$ returns \texttt{unknown} instead of the narrowed type 
\texttt{\{\}} (non-nullish values). This is because the clause diff cannot ``subtract from top''---it 
simply returns the left clause unchanged when $\stypes_1 = \mathsf{None}$.

\end{document}
